\documentclass{article}
\usepackage{graphicx}
\usepackage{algorithm}
\usepackage{algpseudocode}

\title{Formulazione sistema multi-agente a "stato aggregato"}
\author{Lorenzo Pecorari}
\date{November 2025}

\begin{document}
\maketitle

\section{Spazio degli stati e delle osservazioni}
\subsection{Spazio degli stati per l'ambiente}
Siano $I = \{1, 2, ..., n\}$ l'insieme degli indici degli agenti presenti nel sistema e $S$ l'insieme degli stati mantenuto dall'ambiente tale da contenere tutti gli stati $s_t^{env}$ inerenti a tutte le informazioni di ogni agente:\\\\
$s_t^{env} = (b_t^1, \tau_t^1, f_{t-1}^1, x_{t-1}^1, g_{t-1}^1, h_{t-1}^i, ...,
                b_t^i, \tau_t^i, f_{t-1}^i, x_{t-1}^i, g_{t-1}^i, h_{t-1}^i,
                ..., b_t^n, \tau_t^n, f_{t-1}^n, x_{t-1}^n, h_{t-1}^n) = (b_t^i, \tau_t^i, f_{t-1}^i, x_{t-1}^i, g_{t-1}^i, h_{t-1}^i)_{i \in [1, n], t > 0}$ .\\\\
Quindi lo spazio degli stati è definibile come $S = \{s_0\} \cup \{(b_t^i, \tau_t^i, f_{t-1}^i, x_{t-1}^i, g_{t-1}^i, h_{t-1}^i)_{i \in [1, n], t > 0}\}$, con $s_0$ lo stato iniziale e i parametri $b_t^i$ e $\tau_t^i$ la batteria normalizzata e i timestep normalizzati al tempo $t$ da un nodo generico $i$. Le definizione degli altri parametri seguono nei successivi paragrafi.

\subsection{Osservazioni per gli agenti}
Ogni agente non può ottenere tutto lo stato dell'ambiente ad ogni osservazione ma solo una visione parziale che gli permette di avere informazioni sulle proprie specifiche locali come livello di batteria e timestep completati e dettagli a livello globale attraverso variabili aggregate e medie. Ogni osservazione che riceve l'agente $i$ al tempo $t$ è definibile come $o_t^i$ ma per semplicità di notazione viene assunta come "stato" $s_t^i$. Di seguito i dettagli di ognuna di esse:

\begin{itemize}
    \item $s_t^i = (l_t^i, m_t^i)$, stato (osservazione) ottenibile dal nodo $i$ al tempo $t$ dall'ambiente
    \item $l_t^i = (b_t^i, \tau_t^i)$, stato locale del nodo $i$ al tempo $t$ 
    \item $m_t^i = (mb_t^i, m\tau_t^i, mf_t^i, mh_t^i, req_t^i, acp_t^i)$, stato aggregati al tempo t per nodi diversi da $i$
    
    \item $b_t^i \in [0.0, 1.0]$ valore normalizzato di batteria residua di $i$ in $t$
    \item $\tau_t^i \in [0.0, 1.0]$, valore normalizzati di timestep completati da $i$ al tempo $t$ rispetto a $T$
    
    \item $mb_t^i \in [0.0, 1.0]$, media dei valori delle batterie dei nodi $j \neq i$ in $t$
    \item $m\tau_t^i \in [0.0, 1.0]$, media dei timestamp normalizzati e completati dai nodi $j \neq i$ al tempo $t$
    \item $mf_t^i \in \mathbf{N}$, media del framerate locale adottato dai nodi $j \neq i$ in $t-1$
    \item $mh_t^i \in \mathbf{N}$, media del framerate di offloading adottato dai nodi $j \neq i$ al tempo $t-1$
    \item $req_t^i \in \mathbf{N}_{\geq0}$, quantità di nodi diversi da $i$ che in $t-1$ hanno richiesto offloading 
    \item $acp_t^i \in \mathbf{N}_{\geq0}$, quantità di nodi diversi da $i$ che in $t-1$ hanno accettato offloading 
\end{itemize}

\noindent Per ogni agente è possibile definire uno spazio delle osservazioni $S^i \equiv O^i = \{s_0^i, s_1^i, ..., s_j^i, ..., s_k^i\}$ dove ogni stato $j-esimo$ è definito come $s_j^i = (l_j^i, m_j^i)$ .\\

\section{Spazio delle azioni}
Sia A l'insieme dello spazio delle azioni congiunte dell'ambiente:
\begin{itemize}
    \item $A = \{a_1, a_2, ..., a_i, ..., a_k\}$
    \item $a_i = (a_i^1, a_i^2, ..., a_i^j, ..., a_i^n)$
    \item $a_i^j = (f_i^j, x_i^j, g_i^j, h_i^j)$ dove:
    \begin{itemize}
        \item $f_i^j \in [0, fps_{max}^j]$, framerate da adottare localmente per il nodo $j$ con valore massimo $fps_{max}^j$ definito a priori\\
        \item $x_i^j = $\{ \begin{array}{lr}
            $0 , ricezione o invio non abilitati$ \\
            $1 , invio lavoro ad altro nodo$\\
            $2 , ricezione lavoro da altro nodo$
            \end{array}\\
        \item $g_i^j = $\{ \begin{array}{lr}
            $0$, $\space se \space x_i^j$ = 0 \\
            $k \neq j$, se $x_i^j$ = 1 $\\
            \end{array}\\
        \item $h_i^j \in [0, fps_{max}^j]$, framerate per computazione di offloading
    \end{itemize}
\end{itemize}

\noindent Per ogni agente è possibile definire uno spazio delle azioni equivalente a $A^i = \{a_0^i, a_1^i, ..., a_j^i, ..., a_k^i\}$ in cui ogni azione $j-esima$ del nodo $i$ è $a_i^j = (f_i^j, x_i^j, g_i^j, h_i^j)$.

\subsection{Definizione di variabili dello stato aggregato}
Sulla base delle precedenti definizioni, è possibile definire le variabili aggregate prima citate per rendere l'agente informato riguardo all'evoluzione dello stato degli agenti di cui non ha l'informazione completa:

\begin{itemize}

    \item $mb_t^i = \frac{1}{n-1} \sum_{j \neq i}^n b_{t-1}^j$
    \item $m\tau_t^i = \frac{1}{n-1} \sum_{j \neq i}^n \tau_{t-1}^j$
    \item $mf_t^i = \frac{1}{n-1} \sum_{j \neq i}^n f_{t-1}^j$
    \item $mh_t^i = \frac{1}{n-1} \sum_{j \neq i}^n h_{t-1}^j$
    \item $req_t^i = \sum_{j \neq i}^{n} \mathbf{1}(x_{t-1}^j = 1)$
    \item $acp_t^i = \sum_{j \neq i}^{n} \mathbf{1}(x_{t-1}^j = 2)$
     
\end{itemize}

dove la funzione $\mathbf{1}(arg)$ restituisce $1$ se $arg$ risulta vero, $0$ altrimenti.

\section{Nodo - agente}
\subsection{Parametri del nodo}
Definizione di parametri dei nodi:
\begin{itemize}
    \item $B$, capacità massima della batteria in Wh
    \item $P_I$, potenza assorbita in idle in W
    \item $P_M$, potenza assorbita a pieno carico in W
    \item $P_{PV}$, potenza in W generata dal pannello solare
    \item $A_{PV}$, superificie in $m^2$ del pannello
    \item $\eta_{PV}$, efficienza energetica del pannello
    \item $G(t)$, irradianza solare in $W/m^2$ al tempo t
    \item $G_{max}$, irradianza solare in $W/m^2$ massima
    \item $fps_{max}^i$, framerate massimo adaottabile dal nodo
    \item $\Delta t$, intervallo in s con cui simulare la computazione
    \item $E_{F, loc}$, energia in J consumata da per computare una frame localmente
    \item $E_{F, send}$, energia in J consumata per inviare frame ad un nodo per offloading
    \item $E_{F, recv}$, energia in J consumata per ricevere e computare frame da altro nodo
    \item $E_{B, t}$, energia in J della batteria al tempo t
    \item $E_{PV, t}$, energia in J proveniente dal pannello solare al tempo t
\end{itemize}

\subsubsection{Energia dal pannello}
Per ogni timestep $t \in [0, T]$, esiste un valore di $G(t) \in [0, G_{maz}]$ tale per cui si può definire $E_{PV, t} = \int{G(t)\cdot A_{PV} \cdot \eta_{PV} \space dt} \approx {G(t)\cdot A_{PV} \cdot \eta_{PV} \cdot \Delta t $, con $A_{PV} > 0$ e $\eta_{PV} \in[0.0, 1.0]$.

\subsubsection{Energia della batteria}
Per ogni timestep $t \in [0, T]$, la carica della batteria dipende dalla computazione effettuata in $t-1$. Ad ogni esecuzione di un'azione $a_t^i$, il relativo valore dell'energia residua nella batteria sarà determinato da\\\\
$E_{B, t}^i = (E_{B,t-1}^i + E_{PV,t}^i) - ((f_t^i \cdot E_{F, loc}^i \cdot \Delta t)+ (h_t^i \cdot E_{F, send/recv}^i \cdot \Delta t)) $

\subsubsection{Batteria e tempo normalizzati}
Per ogni timestep $t \in [0, T]$, la carica normalizzata della batteria $b_t^i$e i timestep completati $\tau_t^i$ normalizzati rispetto a T sono espressi da:
\begin{itemize}
    \item $b_t^i = \frac{E_{B,t}^i}{B^i}, B^i>0$
    \item $\tau_t^i = \frac{t}{T}, T > 0$
\end{itemize}

\subsection{Tabella Q}
La necessità di rendere l'agente informato sullo stato degli altri nodi senza cadere in una dimensione esponenziale della tabella, come definito in precedenza, porta a considerare valori aggregati e inerenti ai valori medi dei livelli di batteria, dei framerate locali, dei framerate di offloading e il numero di agenti che richiedono e che accettano offloading. Tale approccio permette di mantenere contenuta la crescita del numero di elementi presenti nella Q table. A tal proposito, è possibile considerare lo spazio degli stati $S^i$ come l'insieme degli stati tali che per ogni agente e per ogni timestep si ha $s_t^i = (l_t^i, m_t^i)$. La tabella in questione per ogni agente, quindi, sarà ottenibile dal prodotto cartesiano tra lo spazio degli stati e lo spazio delle azioni: $Q^i =S^i \times A^i $.\\

Localmente, $l_t^i$ definisce il livello della batteria e la percentuale di timestep completati, con 11 valori possibili per entrambi. \\

Gloablmente, in $m_t^i$ si hanno $mf_t^i$ e $mh_t^i$, i valori medi dei framerate adottati localmente da ogni nodi $j \neq i$. Questi hanno come limite superiore $F_t^i = max(fps_{max}^j)_{j\neq i}$ e $H_t^i = max(fps_{max}^j)_{j \neq i}$. Sulla base di tale ragionamento, si assume che i valori considerabili di $mf_t^i$ e $mh_t^i$ saranno rispettivamente $F_t^i+1$ e $H_t^i+1$. Per quanto riguarda il numero di nodi che hanno richiesto ed accettato offloading, $req_t^i$ e $acp_t^i$, questo sarà compreso tra 0 e n; il limite superiore del loro valore sarà quindi $n+1$ .Infine, per i valori medi delle batterie e dei timesteps, il limite superiore sarà uguale a quello delle relative variabili "locali" di $l_t^i$.\\

Si assume quindi che $|S^i| = (11 \cdot 11) \cdot (11 \cdot 11 \cdot F_t^i \cdot H_t^i \cdot (n+1) \cdot (n+1)) = 11^2 \cdot (11^2 \cdot (n+1)^2 \cdot (max(fps_{max}^j)_{j \neq i} +1)^2 = 11^4 \cdot (n+1)^2 \cdot (max(fps_{max}^j)_{j \neq i} +1)^2$ . \\

In maniera simile allo spazio degli stati, è possibile definire lo spazio delle azioni caratterizzato dagli elementi $a_t^i = (f_t^i, x_t^i, g_t^i, h_t^i)$, dove:
\begin{itemize}
    \item $f_t^i \in [0, fps_{max}^i]$, considerabili $fps_{max}^i+1$ valori differenti
    \item $x_t^i \in \{0, 1, 2\}$, considerabili 3 valori differenti
    \item $g_t^i \in [0, n]$, considerabili $n+1$ valori con $n$ il numero di nodi
    \item $h_t^i \in [0, fps_{max}^i]$, considerabili $fps_{max}^i+1$ valori differenti
\end{itemize}

Seguendo l'approccio applicato per lo spazio degli stati, è possibile calcolare la dimensione dello spazio delle azioni "locali" $A^i$: $|A^i| = (fps_{max}^i + 1) \cdot 3 \space \cdot (n+1) \space \cdot (fps_{max}^i + 1) = 3(fps_{max}^i + 1)^2(n+1)$ .\\

Dalla conoscenza di entrambe le dimensioni degli spazi che vengono mantenuti nella Q table da un aegnte $i$, si ricava che $|Q^i| = |S^i \times A^i| = |S^i| \cdot |A^i| = 11^4 \cdot (n+1)^2 \cdot (max(fps_{max}^j)_{j \neq i} +1)^2 \cdot 3(fps_{max}^i + 1)^2(n+1) = 3 \cdot 11^4 \cdot (max(fps_{max}^j)_{j \neq i} +1)^2 \cdot (fps_{max}^i + 1)^2 \cdot (n+1)^3$

\subsubsection{Scelta dell'azione}
Seguendo l'algoritmo $\epsilon-greedy$, si cerca di ottenere un trade-off tra exploration ed exploitation per l'agente, basandosi sul decay factor $\epsilon_{dec} \in[0, 1]$, sul valore iniziale $\epsilon_0$ e su quello finale $\epsilon_{fin}$. Per ogni timestep $t$, l'agente sceglie casualmente l'azione con probabilità $\epsilon$, mentre con probabilità $1 - \epsilon$ "effettua un ragionamento" in modo da ottimizzare la ricompensa.

\subsubsection{Osservazione dell'environment}
Ogni volta che l'agente osserva l'ambiente, l'informazione ritornata è parziale rispetto allo stato di cui tiene traccia l'ambiente ma permette comunque al nodo di avere la conoscenza dei suoi parametri locali e una serie di valori che di stimare la situazione degli altri nodi nel timestep $t$. Rispetto ad avere tutto lo stato congiunto, tale apaproccio permette di mantenere una conoscenza di tutti i nodi senza eccedere con la dimensione della tabella Q.

\subsubsection{Aggiornamento della tabella}
Ad ogni timestep $t$, l'agente sceglie un'azione $a_t^i$ sulla base dell'osservazione ottenuta dall'ambiente. Secondo la policy $\pi_{\theta}^i$, viene scelta un'azione $a_t^i = \pi_\theta^i(s_t)$, eseguita e collezionato il reward. La policy da seguire, in questo caso, prevede la scelta  dell'azione che massimizza il valore Q in corrispondenza di tale stato $s_t$: $a_t^i = argmax_{a \in A^i} Q(s_t^i, a)$.\\\\ A seguito di ciò si procede con l'aggiornamento della tabella secondo:\\

$Q(s_t^i, a_t^i) \leftarrow (1-\alpha)Q(s_t^i, a_t^i) + \alpha(r_t^i + \gamma(max_{a_{t+1}^i} Q(s_{t+1}^i, a_{t+1}^i)))$ \\\\
con $\alpha \in [0, 1]$ come learning rate e $\gamma \in [0, 1]$ come discount factor. 


\section{Funzione di reward}
Dati due nodi $i \in I$ e $j \in I$ tali che $j \neq i$, per ogni timestep $t \in [0, T]$, con azioni $a_t^i = (f_t^i, x_t^i, g_t^i, h_t^i)$ e $a_t^j = (f_t^j, x_t^j, g_t^j, h_t^j)$, è possibile definire la funzione di reward che permetta ad essi di ottenere una ricompensa sulla base dello stato corrente, dell'azione eseguita e dello stato successivo:\\

$R_i(t) = r_t^i = r_{frames, t}^i + r_{battery, t}^i + r_{cooperation, t}^i$ 
\begin{algorithm}
    \State{// Verifica se la somma degli fps adottati localmente e quelli di offloading sono minori uguali a quelli massimi per il nodo}
    \caption{Algoritmo per ricompensa per gestione delle frame, $r_{frames, t}^i$}\label{alg:cap}
    \begin{algoerithmic}
    \If{($(f_t^i + h_t^i) \leq fps_{max}^i$)}
        \State{return 1}
        %\State{return $\frac{(f_t^i + h_t^i)}{fps_{max}^i}$}
    \Else
    \State{return 0}
    \EndIf
    \end{algoerithmic}
\end{algorithm}

\begin{algorithm}
    \caption{Algoritmo per ricompensa per gestione della batteria, $r_{battery, t}^i$}\label{alg:cap}
    \begin{algoerithmic}
    \State{// Verifica se $i$ non effettua offloading}
    \If{($x_t^i =0$ \land $((f_t^i \cdot E_{F, loc} \cdot \Delta t) < (E_{B,t}^i + E_{PV, t}^i))$)}
        \State{return 1}\\
        %\State{return $(f_t^i \cdot E_{F, loc} \cdot \Delta t)$/ $(E_{B,t}^i + E_{PV, t}^i)$}\\

    \State{// Se $i$ invia lavoro, l'energia a disposizione deve essere necessaria}
    \ElsIf{($x_t^i =1$ \land $(((f_t^i \cdot E_{F, loc} \cdot \Delta t) + (h_t^i \cdot E_{F, send} \cdot \Delta t))) < (E_{B,t}^i + E_{PV, t}^i))$)}
        \State{return 1}\\
        
    %\State{return $((f_t^i \cdot E_{F, loc} \cdot \Delta t) + (h_t^i \cdot E_{F, send} \cdot \Delta t)) / (E_{B,t}^i + E_{PV, t}^i)$}\\


    \State{// Se $i$ riceve lavoro, l'energia a disposizione deve essere necessaria}
    \ElsIf{($x_t^i =2$ \land $(((f_t^i \cdot E_{F, loc} \cdot \Delta t) + (h_t^i \cdot E_{F, recv} \cdot \Delta t))) < (E_{B,t}^i + E_{PV, t}^i))$)}
        \State{return 1}\\
        %\State{return $((f_t^i \cdot E_{F, loc} \cdot \Delta t) + (h_t^i \cdot E_{F, recv} \cdot \Delta t)) / (E_{B,t}^i + E_{PV, t}^i)$}\\

    \State{// Vincolo non soddisfatto, ricompensa nulla}
    \Else
    \State{return 0}
    \EndIf
    \end{algoerithmic}
\end{algorithm}

\begin{algorithm}
    \caption{Algoritmo per ricompensa per cooperazione tra i nodi $r_{cooperation, t}^i$}\label{alg:cap}
    \begin{algoerithmic}

    \State{// Se $i$ non effettua offloading e non designa nessuno per tale operazione}
    \If{($x_t^i =0$ \land $g_t^i = 0$)}
        \State{return 1}\\

    \State{// Se $i$ effettua offloading, invia lavoro, designa un altro nodo e tale nodo intende ricevere da i}
    \ElsIf{($x_t^i =1$ \land $g_t^i \neq 0 \land g_t^i \neq i \land x_t^{g_t^i} = 2 \land g_t^{g_t^i} = i$)}
    \State{return 1}\\

    \State{// Se $i$ effettua offloading, riceve lavoro, designa un altro nodo e tale nodo intende inviare a i}
    \ElsIf{(($x_t^i =2$ \land $g_t^i \neq 0 \land g_t^i \neq i \land x_t^{g_t^i}= 1 \land g_t^{g_t^i} = i$))}
    \State{return 1}\\

    \State{// Nessun vincolo soddisfatto, ricompensa nulla}
    \Else
    \State{return 0}
    \EndIf
    \end{algoerithmic}
\end{algorithm}


\end{document}
